\section{Exercices d'application}

\rubrique{Principe additif et multiplicatif}

\exo{}\diff{1}
Un restaurant propose un menu composé d'une entrée, d'un plat principal et d'un dessert. Il y a 3 entrées, 5 plats principaux et 4 desserts. Combien de menus différents peut-on composer ?

\exo{}\diff{1}
Dans une classe de Première C, on doit élire un délégué et un suppléant parmi 12 candidats. Combien de couples (délégué, suppléant) distincts peut-on former ?

\exo{}\diff{2}
Un code secret est formé de 4 chiffres (de 0 à 9) suivis de 2 lettres (de A à Z). Combien de codes différents peut-on créer ?

\rubrique{p-listes et arrangements}

\exo{}\diff{1}
On lance 3 fois de suite un dé équilibré à 6 faces. Combien de résultats différents peut-on obtenir ?

\exo{}\diff{2}
Une urne contient 10 boules numérotées de 1 à 10. On tire successivement et sans remise 3 boules. Combien de tirages ordonnés différents peut-on obtenir ?

\exo{}\diff{2}
De combien de façons peut-on asseoir 5 personnes sur 7 chaises alignées ?

\rubrique{Permutations}

\exo{}\diff{1}
Combien de mots de 5 lettres distinctes peut-on former avec les lettres du mot \textit{CHIEN} ?

\exo{}\diff{2}
On dispose de 6 livres différents que l'on veut ranger sur une étagère. Combien de rangements différents peut-on faire ?

\exo{}\diff{2}
Combien d'anagrammes du mot \textit{BAC} peut-on écrire ?

\rubrique{Combinaisons}

\exo{}\diff{1}
Dans une classe de 20 élèves, on doit choisir un groupe de 3 élèves pour un projet. Combien de groupes différents peut-on former ?

\exo{}\diff{2}
Une main de 5 cartes est tirée d'un jeu de 32 cartes. Combien de mains différentes contiennent exactement 3 rois ?

\exo{}\diff{3}
Une commission de 4 personnes doit être formée parmi 6 hommes et 5 femmes. Combien de commissions peut-on former si elle doit comporter au moins 2 femmes ?

\rubrique{Propriétés des combinaisons et triangle de Pascal}

\exo{}\diff{2}
Calculer $\binom{7}{3}$ et $\binom{7}{4}$. Que remarque-t-on ?

\exo{}\diff{2}
En utilisant la relation de Pascal, calculer $\binom{8}{3}$ sachant que $\binom{7}{2}=21$ et $\binom{7}{3}=35$.

\exo{}\diff{3}
Démontrer que $\binom{n}{k} = \binom{n}{n-k}$ pour $0 \le k \le n$.

\rubrique{Binôme de Newton}

\exo{}\diff{2}
Développer $(x+1)^4$ à l'aide du binôme de Newton.

\exo{}\diff{3}
Déterminer le coefficient de $a^3b^2$ dans le développement de $(a+b)^5$.

\exo{}\diff{3}
En utilisant le binôme de Newton, calculer la somme $\sum_{k=0}^{n} \binom{n}{k}$.

\section{Exercices d'entraînement}

\rubrique{Principe additif et multiplicatif}

\exo{}\diff{1}
Un magasin propose 4 modèles de chemises, 3 modèles de pantalons et 2 modèles de vestes. Combien de tenues différentes (chemise + pantalon + veste) peut-on composer ?

\exo{}\diff{1}
Pour se rendre de Douala à Yaoundé, on peut prendre 3 compagnies de bus, 2 compagnies de train ou 1 compagnie d'avion. Combien de choix de transport a-t-on ?

\exo{}\diff{2}
Un code de carte bancaire est composé de 4 chiffres. Combien de codes différents existe-t-il ?

\exo{}\diff{2}
Dans un restaurant, on peut choisir entre 2 entrées, 3 plats et 2 desserts. Combien de repas complets différents peut-on commander ?

\rubrique{p-listes et arrangements}

\exo{}\diff{1}
On lance une pièce de monnaie 4 fois de suite. Combien de suites de résultats (Pile ou Face) peut-on obtenir ?

\exo{}\diff{2}
Une urne contient 8 boules numérotées de 1 à 8. On tire successivement et sans remise 4 boules. Combien de tirages ordonnés différents ?

\exo{}\diff{2}
De combien de façons peut-on placer 3 personnes sur 5 chaises ?

\exo{}\diff{3}
Combien de mots de 4 lettres (avec répétition possible) peut-on former avec les lettres A, B, C ?

\rubrique{Permutations}

\exo{}\diff{1}
Combien de mots de 4 lettres distinctes peut-on former avec les lettres du mot \textit{MAIL} ?

\exo{}\diff{2}
Combien d'anagrammes du mot \textit{MATHS} peut-on écrire ?

\exo{}\diff{2}
On dispose de 8 livres différents. Combien de rangements différents peut-on faire sur une étagère ?

\exo{}\diff{3}
Combien d'anagrammes du mot \textit{ANANAS} peut-on écrire ?

\rubrique{Combinaisons}

\exo{}\diff{1}
Dans un groupe de 15 personnes, on choisit 2 personnes pour former un binôme. Combien de binômes différents ?

\exo{}\diff{2}
Un jeu de 52 cartes. Combien de mains de 5 cartes contiennent exactement 2 as ?

\exo{}\diff{2}
Une équipe de 3 garçons et 2 filles doit être formée parmi 8 garçons et 6 filles. Combien d'équipes différentes ?

\exo{}\diff{3}
Un comité de 5 personnes est choisi parmi 7 hommes et 4 femmes. Combien de comités comportant au moins 3 hommes ?

\rubrique{Propriétés des combinaisons et triangle de Pascal}

\exo{}\diff{2}
Calculer $\binom{6}{2}$ et $\binom{6}{4}$. Que remarque-t-on ?

\exo{}\diff{2}
En utilisant la relation de Pascal, calculer $\binom{9}{4}$ sachant que $\binom{8}{3}=56$ et $\binom{8}{4}=70$.

\exo{}\diff{3}
Démontrer la relation de Pascal : $\binom{n}{k} = \binom{n-1}{k-1} + \binom{n-1}{k}$.

\rubrique{Binôme de Newton}

\exo{}\diff{2}
Développer $(2x+1)^3$ à l'aide du binôme de Newton.

\exo{}\diff{3}
Déterminer le coefficient de $x^2y^3$ dans le développement de $(x+y)^5$.

\exo{}\diff{3}
Calculer $\sum_{k=0}^{5} \binom{5}{k} 2^k$.

\section{Problèmes type Probatoire/Bac}

\sujetbac[Probatoire Blanc — Dénombrement]
Un sac contient 12 billes : 5 rouges, 4 vertes et 3 bleues. On tire simultanément 4 billes du sac.
\begin{enumerate}[label=\arabic*)]
\item Combien de tirages différents peut-on effectuer ?
\item Combien de tirages contiennent exactement 2 billes rouges ?
\item Combien de tirages contiennent au moins une bille verte ?
\item Combien de tirages contiennent des billes de toutes les couleurs ?
\end{enumerate}

\sujetbac[Baccalauréat Blanc — Anagrammes et combinaisons]
On considère le mot \textit{MATHÉMATIQUE}.
\begin{enumerate}[label=\arabic*)]
\item Combien d'anagrammes de ce mot peut-on former ?
\item Combien d'anagrammes commencent par la lettre M ?
\item Combien d'anagrammes ont toutes les lettres A regroupées ?
\item On choisit au hasard 3 lettres parmi les lettres du mot. Combien de choix différents peut-on faire ?
\end{enumerate}

\sujetbac[Probatoire — Comité et binôme]
Une association compte 8 hommes et 7 femmes. On doit former un bureau de 4 personnes.
\begin{enumerate}[label=\arabic*)]
\item Combien de bureaux différents peut-on former ?
\item Combien de bureaux comportent exactement 2 hommes ?
\item Combien de bureaux comportent au moins une femme ?
\item On décide que le bureau doit comporter un président, un secrétaire et deux membres. Combien de bureaux distincts peut-on former ?
\end{enumerate}

\section{Corrigés}

\corrige{de l'exercice 1}
Par le principe multiplicatif, le nombre de menus est $3 \times 5 \times 4 = 60$.

\corrige{de l'exercice 2}
Pour le délégué, on a 12 choix. Pour le suppléant, il reste 11 choix. Soit $12 \times 11 = 132$ couples.

\corrige{de l'exercice 3}
Pour les 4 chiffres : $10^4 = 10\,000$ possibilités. Pour les 2 lettres : $26^2 = 676$ possibilités. Par le principe multiplicatif, $10\,000 \times 676 = 6\,760\,000$ codes.

\corrige{de l'exercice 4}
Chaque lancer a 6 issues. On a une 3-liste avec répétition : $6^3 = 216$ résultats.

\corrigé{de l'exercice 5}
Il s'agit d'un arrangement sans répétition de 3 boules parmi 10 : $A_{10}^3 = 10 \times 9 \times 8 = 720$.

\corrigé{de l'exercice 6}
On choisit 5 chaises parmi 7, puis on permute les 5 personnes sur ces chaises. Soit $\binom{7}{5} \times 5! = 21 \times 120 = 2\,520$. Ou directement $A_7^5 = 7 \times 6 \times 5 \times 4 \times 3 = 2\,520$.

\corrigé{de l'exercice 7}
Il s'agit d'une permutation de 5 lettres distinctes : $5! = 120$ mots.

\corrigé{de l'exercice 8}
Permutation de 6 livres distincts : $6! = 720$ rangements.

\corrigé{de l'exercice 9}
Le mot \textit{BAC} a 3 lettres distinctes. Nombre d'anagrammes : $3! = 6$.

\corrigé{de l'exercice 10}
Nombre de combinaisons de 3 élèves parmi 20 : $\binom{20}{3} = \frac{20 \times 19 \times 18}{3 \times 2 \times 1} = 1\,140$.

\corrigé{de l'exercice 11}
Il y a 4 rois dans le jeu. On choisit 3 rois parmi 4 : $\binom{4}{3} = 4$. On choisit les 2 autres cartes parmi les 28 non-rois : $\binom{28}{2} = 378$. Soit $4 \times 378 = 1\,512$ mains.

\corrigé{de l'exercice 12}
Cas 1 : 2 femmes et 2 hommes : $\binom{5}{2} \times \binom{6}{2} = 10 \times 15 = 150$.
Cas 2 : 3 femmes et 1 homme : $\binom{5}{3} \times \binom{6}{1} = 10 \times 6 = 60$.
Cas 3 : 4 femmes : $\binom{5}{4} \times \binom{6}{0} = 5 \times 1 = 5$.
Total : $150 + 60 + 5 = 215$ commissions.

\corrigé{de l'exercice 13}
$\binom{7}{3} = \frac{7 \times 6 \times 5}{3 \times 2 \times 1} = 35$ et $\binom{7}{4} = \frac{7 \times 6 \times 5 \times 4}{4 \times 3 \times 2 \times 1} = 35$. On remarque que $\binom{7}{3} = \binom{7}{4}$.

\corrigé{de l'exercice 14}
Par la relation de Pascal : $\binom{8}{3} = \binom{7}{2} + \binom{7}{3} = 21 + 35 = 56$.

\corrigé{de l'exercice 15}
$\binom{n}{k} = \frac{n!}{k!(n-k)!}$ et $\binom{n}{n-k} = \frac{n!}{(n-k)!(n-(n-k))!} = \frac{n!}{(n-k)!k!}$. Les deux expressions sont égales.

\corrigé{de l'exercice 16}
$(x+1)^4 = \sum_{k=0}^{4} \binom{4}{k} x^{4-k} 1^k = \binom{4}{0}x^4 + \binom{4}{1}x^3 + \binom{4}{2}x^2 + \binom{4}{3}x + \binom{4}{4} = x^4 + 4x^3 + 6x^2 + 4x + 1$.

\corrigé{de l'exercice 17}
Dans $(a+b)^5$, le terme général est $\binom{5}{k} a^{5-k} b^k$. Pour $a^3b^2$, on a $5-k=3$ et $k=2$. Le coefficient est $\binom{5}{2} = 10$.

\corrigé{de l'exercice 18}
$\sum_{k=0}^{n} \binom{n}{k} = (1+1)^n = 2^n$ d'après le binôme de Newton.

\corrigé{du problème 1}
\begin{enumerate}[label=\arabic*)]
\item Nombre total de tirages de 4 billes parmi 12 : $\binom{12}{4} = 495$.
\item Exactement 2 rouges : $\binom{5}{2} \times \binom{7}{2} = 10 \times 21 = 210$.
\item Au moins une verte : on calcule le complément (aucune verte). Aucune verte : on tire parmi les 8 non-vertes : $\binom{8}{4} = 70$. Donc $495 - 70 = 425$.
\item Toutes les couleurs : on a 1 rouge, 1 verte, 1 bleue et une quatrième bille d'une couleur quelconque. Cas 1 : 2 rouges, 1 verte, 1 bleue : $\binom{5}{2} \times 4 \times 3 = 10 \times 4 \times 3 = 120$. Cas 2 : 1 rouge, 2 vertes, 1 bleue : $5 \times \binom{4}{2} \times 3 = 5 \times 6 \times 3 = 90$. Cas 3 : 1 rouge, 1 verte, 2 bleues : $5 \times 4 \times \binom{3}{2} = 5 \times 4 \times 3 = 60$. Total : $120 + 90 + 60 = 270$.
\end{enumerate}

\corrigé{du problème 2}
Le mot \textit{MATHÉMATIQUE} a 12 lettres : M (2 fois), A (2 fois), T (2 fois), H (1 fois), É (1 fois), I (1 fois), Q (1 fois), U (1 fois), E (1 fois).
\begin{enumerate}[label=\arabic*)]
\item Nombre d'anagrammes : $\frac{12!}{2!\,2!\,2!} = \frac{479\,001\,600}{8} = 59\,875\,200$.
\item Commençant par M : on fixe un M, il reste 11 lettres avec M (1 fois), A (2 fois), T (2 fois) : $\frac{11!}{2!\,2!} = \frac{39\,916\,800}{4} = 9\,979\,200$.
\item Toutes les lettres A regroupées : on considère le bloc AA comme une lettre. On a alors 11 éléments (bloc AA + 10 autres lettres) avec M (2 fois), T (2 fois) : $\frac{11!}{2!\,2!} = 9\,979\,200$.
\item On choisit 3 lettres parmi les 9 lettres distinctes (M, A, T, H, É, I, Q, U, E) : $\binom{9}{3} = 84$.
\end{enumerate}

\corrigé{du problème 3}
\begin{enumerate}[label=\arabic*)]
\item Nombre total de bureaux de 4 personnes parmi 15 : $\binom{15}{4} = 1\,365$.
\item Exactement 2 hommes : $\binom{8}{2} \times \binom{7}{2} = 28 \times 21 = 588$.
\item Au moins une femme : complément (aucune femme, soit 4 hommes) : $\binom{8}{4} = 70$. Donc $1\,365 - 70 = 1\,295$.
\item On choisit d'abord les 4 personnes : $\binom{15}{4} = 1\,365$. Ensuite, on attribue les rôles : président (4 choix), secrétaire (3 choix), les deux autres sont membres. Soit $1\,365 \times 4 \times 3 = 16\,380$.
\end{enumerate}